\documentclass{article}
\usepackage{amsmath, amssymb, amsthm}
\usepackage{hyperref}
\usepackage{geometry}
\geometry{margin=1in}

\title{Erd\H{o}s Problem \#1209}
\date{Last updated: 2026-04-04}
\author{Source: erdosproblems.com}

\begin{document}
\maketitle

\noindent\textbf{Status:} OPEN \quad \textbf{No prize}

\noindent\textbf{Tags:} number theory

\medskip

\noindent\textbf{Problem Statement:}

Let $A=\{a_1&#60;a_2&#60;\cdots\}$ be a sequence of integers which tends to infinity sufficiently fast. If there is an $n$ such that all $n+a_k$ are primes then must there exist infinitely many such $n$? What if we ask for $n+a_k$ to be squarefree instead of prime? Are there $n$ such that $n+2^{2^k}$ is always a prime (or always squarefree, or infinitely often a prime, or infinitely often squarefree)?




For the first two questions Erd\H{o}s \cite{Er80} wrote 'unless I overlook a trivial way of getting a counterexample these questions are quite hopeless'. There is indeed a trivial counterexample (a variant of the construction in [429] ): define $a_1=2$ and for $k\geq 2$ let $a_k&gt;a_{k-1}$ be a prime such that $a_k+k\equiv 0\pmod{q_k}$, where $q_k$ is some prime not dividing $k$. This sequence can be made to grow arbitrarily fast. A similar construction with $\pmod{q_k^2}$ provides a counterexample to the squarefree question. ebarschkis and GPT have proved that there are no $n$ such that $n+2^{2^k}$ is always prime: let $n\geq 3$ be any odd integer. If $k$ is chosen sufficiently large, and $p=n+2^{2^{k}}$ is prime, then the multiplicative order of $2^{2^k}\pmod{p}$, say $m$ is odd, and hence if $l$ is chosen such that $2^l\equiv 1\pmod{m}$ then $p\mid n+2^{2^{k+rl}}$ for all $r\geq 1$. See also [429] and [1102] .




References

[Er80] Erd\H{o}s, Paul, A survey of problems in combinatorial number theory . Ann. Discrete Math. (1980), 89-115.

\noindent\textbf{Related OEIS sequences:} possible


\bigskip
\noindent\small{Source: \url{https://www.erdosproblems.com/1209}}
\end{document}
